\documentclass[11pt]{article}
\usepackage[T1]{fontenc}
\usepackage{lmodern,amsmath,amssymb,amsthm,booktabs,geometry,microtype,xcolor,hyperref}
\geometry{margin=1.03in}
\definecolor{ink}{HTML}{15475C}
\hypersetup{colorlinks=true,linkcolor=ink,urlcolor=ink,citecolor=ink,pdftitle={A full-dimensional null set of uniquely optimal Egyptian fractions}}
\newtheorem{theorem}{Theorem}[section]\newtheorem{lemma}[theorem]{Lemma}\newtheorem{corollary}[theorem]{Corollary}\newtheorem{proposition}[theorem]{Proposition}
\newcommand{\HH}{\dim_{\mathrm H}}\newcommand{\PP}{\dim_{\mathrm P}}\newcommand{\U}{\mathcal U}\newcommand{\N}{\mathbb N}
\title{A full-dimensional null set of uniquely optimal Egyptian fractions}
\author{Research note}\date{27 November 2026}
\begin{document}\maketitle
\begin{abstract}\noindent
For which real numbers does choosing the largest admissible unit fraction at each step give the unique best approximation at every order? We prove that these numbers have Hausdorff dimension one, even when competitors may repeat denominators. Kova\v c's theorem implies that they have Lebesgue measure zero. We determine the joint Hausdorff and packing spectrum of the lower and upper double-exponential rates of the optimal errors inside this set. The main estimate counts competing sorted lists according to their first departure from the intended list. Separation at every earlier partial sum forces a bound on each subsequent rival denominator, and the resulting losses are summable. The proof is unconditional.
\end{abstract}

\section{The exceptional set}
For $0<x<1$, start at $s_0=0$ and put
\[
 a_{n+1}=\lfloor (x-s_n)^{-1}\rfloor+1,\qquad s_{n+1}=s_n+1/a_{n+1}.
\]
The inequality $s_{n+1}<x$ is strict. The greedy denominators increase strictly since $x-s_{n+1}\le 1/[a_{n+1}(a_{n+1}-1)]$. Write $E_n=x-s_n$. Call $x$ \emph{uniquely optimal at every order} if, for each $n$, every nondecreasing list $b_1,\ldots,b_n$ of positive integers with $\sum 1/b_i<x$ has sum at most $s_n$, with equality only for the greedy list. Let $\U$ denote the set of these $x$. Allowing repeated entries makes this the stronger of two common optimality conventions.

\begin{theorem}\label{thm:full}
The set $\U$ has Lebesgue measure zero and Hausdorff dimension one. Its transcendental elements alone have Hausdorff dimension one. The same conclusions hold when competitors must have distinct denominators.
\end{theorem}

The measure assertion follows from Kova\v c's theorem for distinct denominators \cite{kovac}. Every point of $\U$ belongs to the set in that theorem, which requires only one nested sequence of best approximations from some order onward. Nathanson asked for irrational examples of all-order unique optimality \cite{nathanson}. Kova\v c and Tang constructed such an example and proved eventual greediness for every positive rational \cite{kt}. Theorem~\ref{thm:full} gives the Hausdorff size of the exceptional set. The conjecture concerning eventual greediness of algebraic irrationals remains outside its scope. Wu studied dimensions for growth restrictions on the denominators in greedy Sylvester expansions \cite{wu}. Our construction compares each prefix with every Egyptian representation of the same length.

For a finer statement define
\[
 \alpha(x)=\liminf_{n\to\infty}\frac{\log\log(1/E_n)}n,
 \qquad \beta(x)=\limsup_{n\to\infty}\frac{\log\log(1/E_n)}n.
\]
All logarithms are natural. The recurrence for $a_n$ gives $\alpha(x)\ge\log 2$: indeed $a_{n+1}-1\ge a_n(a_n-1)\ge(a_n-1)^2$, and eventually $a_n>2$. Also $1/a_{n+1}<E_n\le1/(a_{n+1}-1)$.

\begin{theorem}\label{thm:spectrum}
For $\log 2\le A\le B\le\infty$, let $\U_{A,B}=\{x\in\U:\alpha(x)=A,\ \beta(x)=B\}$. Then
\[
  \HH\U_{A,B}=\frac1{e^B-1},\qquad \PP\U_{A,B}=1,
\]
with the first value interpreted as zero for $B=\infty$. Both dimensions are unchanged after deleting all algebraic numbers, and also if uniqueness is required only among distinct-denominator competitors.
\end{theorem}

For example, the best error can have lower exponent $\log 2$ and upper exponent $\log 3$ on a set of Hausdorff dimension $1/2$. Requiring upper exponent $\log 3$ and lower exponent $\log 3$ gives the same dimension. Both sets have packing dimension one. At $A=B=\log 2$ Theorem~\ref{thm:spectrum} gives Theorem~\ref{thm:full}.

The starting point in ProveIt is the description of heat-kernel decay through approximation by short sums of zero angles. Its section on Li kernels uses a quadratic spacing scale and a Cantor construction \cite{repo}. For the alphabet $1/m$ the spacing is exactly $1/[m(m+1)]$. The extra issue is whether the same partial sums win the optimization at \emph{all} orders. Sections~\ref{sec:tree} and \ref{sec:count} supply this step. The companion paper treats the joint decay question for the original zero angles.

\section{Separated prefixes}\label{sec:tree}
Fix any real sequence $v_n\ge\log 2$. Start with $M_1\ge 8$ integral and choose integer scales
\begin{equation}\label{eq:scale}
 M_{n+1}=\left\lceil (128(n+1))^{3(n+1)} M_n^{\exp(v_n)}\right\rceil,
 \qquad D_n=\{M_n,\dots,2M_n-1\}.
\end{equation}
Independently choose each $a_i$ uniformly from $D_i$. For a prefix let $s_j=\sum_{i\le j}1/a_i$, and let $x=\sum_{i\ge1}1/a_i$. Denote the product probability by $\mathbb P$ and its image on the line by $\mu$. Write
\[
 h_j=\frac1{32M_j^2},\qquad P_j=\prod_{i=1}^j M_i,\qquad P_0=1.
\]
The large prefactor in \eqref{eq:scale} makes elementary numerical inequalities uniform in $j$. The logarithmic rates of interest are unaffected by this factor. The tail satisfies
\begin{equation}\label{eq:tail}
 \frac{1}{2M_{j+1}}< x-s_j \le \frac2{M_{j+1}}<\frac{h_j}{4}.
\end{equation}
Consequently every branch has its displayed denominators as its greedy expansion: $x-s_{j-1}$ is strictly between $1/a_j$ and $1/(a_j-1)$. In particular the coding is injective.

Call a prefix of length $j$ \emph{isolated} if every other sorted list of at most $j$ positive denominators has sum at distance greater than $h_j$ from $s_j$. Distinct lists with the very same sum count as competitors. The empty list is allowed. Isolation at every order implies $x\in\U$ by \eqref{eq:tail}. At the first level every prefix is isolated, by the spacing between consecutive reciprocals.

For later counting, centers corresponding to different length-$j$ words are separated by at least $1/(8M_j^2)$. In fact, if the first difference is at $i$, the two reciprocal terms differ by at least $1/(4M_i^2)$, whereas the two nonnegative tails each lie in $[0,2/M_{i+1}]$. The same estimate applies to two infinite branches with this first difference. In particular a single real number can lie within $h_j$ of at most one candidate center of level $j$.

\section{Counting the competitors}\label{sec:count}
\begin{proposition}\label{prop:survive}
With the scales \eqref{eq:scale}, the probability that all prefixes are isolated exceeds $1/2$, for every choice of the sequence $(v_n)$.
\end{proposition}
\begin{proof}
Consider a word that first fails isolation at level $j\ge2$. Any threatening list of length at most $j-1$ would have sum within $1/M_j+h_j<h_{j-1}$ of $s_{j-1}$. Isolation forces it to be the exact parent list, whose sum is at distance $1/a_j>h_j$ from $s_j$. Thus the threatening list has exactly $j$ entries. Write its sorted entries as $b_1,\ldots,b_j$ and its sum as $z$. Its first disagreement with the word occurs at some $l<j$, since changing only the last denominator changes the sum by at least $1/[a_j(a_j+1)]>h_j$.

The rival tail from $l$ onward sums to at least $1/a_l-h_j\ge 1/(3M_l)$, so $b_l\le 3jM_l$. For each $i>l$ the rival prefix $t_{i-1}=\sum_{k<i}1/b_k$ differs as a list from the chosen prefix. Isolation gives $|t_{i-1}-s_{i-1}|>h_{i-1}$. Moreover
\[
 z\le s_{i-1}+2/M_i+h_j<s_{i-1}+h_{i-1}/2,
 \qquad z\ge s_{i-1}-h_j.
\]
Since $t_{i-1}\le z$, the only possible side is $t_{i-1}<s_{i-1}-h_{i-1}$. The remaining terms must supply at least $h_{i-1}/2$. They are each at most $1/b_i$, hence
\begin{equation}\label{eq:cost}
 b_i\le 2j/h_{i-1}=64jM_{i-1}^2\qquad(l<i\le j).
\end{equation}
This is the gain over counting arbitrary Egyptian representations. After the first disagreement, each successive rival denominator is controlled by the preceding scale.

For fixed $l$ the number of possible rival lists is bounded by
\[
 P_{l-1}(3jM_l)\prod_{i=l+1}^j(64jM_{i-1}^2).
\]
We can drop sortedness for this upper bound. Each listed rival spoils at most one length-$j$ word, by the separation of centers. Divide by $P_j$ and sum over $l$. Since each ratio $M_{i-1}^2/M_i\le1$, the fraction of all words rejected for the first time at level $j$ is at most
\begin{equation}\label{eq:bound}
 j(64j)^j\frac{M_{j-1}^2}{M_j}
 \le \frac{j(64j)^j}{(128j)^{3j}}<2^{-j-1}.
\end{equation}
These are unconditional probabilities on the original tree; no equal-survival assumption about distinct parents enters the estimate. The sum over $j\ge2$ is less than $1/2$.
\end{proof}

The survivor event $S$ is closed in the product space, being the intersection of cylinder conditions. Its image is compact. Restrict $\mu$ to this image to obtain $\nu$, a measure of mass $>1/2$ dominated by $\mu$. It has no atoms because the coding is injective and $P_n\to\infty$. We can henceforth delete any countable set, including all algebraic numbers, without changing its mass.

\section{Dimensions of the surviving branches}
Put $u_n=\log M_n$ and $W_n=\sum_{i=1}^n u_i=\log P_n$. Equation \eqref{eq:scale} gives
\begin{equation}\label{eq:mean}
 u_{n+1}=e^{v_n}u_n+O(n\log(n+2)),\qquad
 \log u_{n+1}=\log u_1+\sum_{i=1}^n v_i+O(1).
\end{equation}
The last error is bounded uniformly in $n$ because $u_n\ge2^{n-1}u_1$. By \eqref{eq:tail}, $\log\log(1/E_n)=\log u_{n+1}+o(1)$ on every branch. The prescribed rates will be the lower and upper limits of the arithmetic means of the $v_i$.

\begin{lemma}\label{lem:dimension}
If $B<\infty$ and $\limsup v_n\le B$, every set of positive $\nu$-measure has Hausdorff dimension at least $1/(e^B-1)$. If there are intervals of indices tending to infinity, of lengths tending to infinity, on which $v_n=\log2$, every set of full $\nu$-measure has packing dimension one.
\end{lemma}
\begin{proof}
For any $Q>e^B$, \eqref{eq:mean} gives $u_{n+1}\le Q u_n$ eventually. Summing backwards gives
\[
 \liminf_{n\to\infty} W_n/u_{n+1}\ge\sum_{i\ge1}Q^{-i}=1/(Q-1).
\]
Thus for every $0<d<1/(e^B-1)$, $P_n^{-1}\le C_d M_{n+1}^{-d}$.

An interval of length $1/M_{n+1}\le r\le1/M_n$ meets a bounded number of distinct parent cylinders of level $n-1$, by the first-difference separation and the tail bound. Level-$n$ cylinders have mass $1/P_n$ and mutual separation at least $1/(8M_n^2)$. It follows that
\[
 \nu(I)\le\mu(I)\le C\min\big\{P_{n-1}^{-1},\,(1+rM_n^2)P_n^{-1}\big\}.
\]
For $r\le M_n^{-2}$ this is $O_d(M_{n+1}^{-d})=O_d(r^d)$. For $M_n^{-2}\le r\le M_n^{-1}$ it is $O_d(rM_n^{1-d})=O_d(r^d)$. The mass distribution principle proves the Hausdorff assertion.

Let $n$ be one past the last index of one of the stated intervals. For every fixed $i\ge1$, $u_{n-i}/u_n\to2^{-i}$ along these endpoints, by \eqref{eq:mean}. As $u_{j+1}\ge2u_j$, it follows that $W_{n-1}/u_n\to1$. Every interval of length $1/M_n$ has mass at most $C/P_{n-1}$. Any set of positive $\nu$ outer measure therefore needs $\gg P_{n-1}$ intervals of this size to cover it, and has upper box dimension one. In a countable cover of a full-measure set at least one member has positive outer measure. The characterization of packing dimension by countable covers and upper box dimension now proves the assertion.
\end{proof}

\section{Choosing the rates and proving sharpness}
To obtain $(A,B)$, use stages $k=1,2,\dots$. Put $A_k=\min(A,\log2+k)$ and $B_k=\min(B,\log2+k)$. Append entries $v_i=A_k$ until the current mean is within $1/k$ of $A_k$ and the current length is at least $k^4$. Append at least one entry $B_k$ and continue at this value until the mean is within $1/k$ of $B_k$. Append $k$ entries equal to $\log2$. During each constant block the mean moves monotonically towards its target. The last $k$ entries change the mean by $O(k^2/k^4)$ since the mean before that block is at most $\log2+k$. Hence the successive means have liminf $A$ and limsup $B$, including when either is infinite. For finite $B$ every entry is at most $B$. Lemma~\ref{lem:dimension} gives the claimed lower dimensions, and gives positive mass and full packing dimension also at $B=\infty$.

For the matching upper bound consider every ordinary greedy expansion, even without optimality. Write $w_i=\log a_i$, $T_n=\sum_{i\le n}w_i$. If $2<R<e^{\beta(x)}$, infinitely often
\[
 w_{n+1}\ge(R-1)T_n.
\]
Otherwise $T_{n+1}<RT_n$ eventually, so $\limsup n^{-1}\log w_{n+1}\le\log R$, contrary to the definition of $\beta(x)$ and the bounds on $E_n$. Each qualifying prefix puts $x$ into an interval to the right of $s_n$ of length at most $2(\prod_{i\le n}a_i)^{-(R-1)}$. For $p=d(R-1)>1$ the $d$th powers of these lengths have finite sum over all increasing lists of all lengths, since
\[
 \sum_{n\ge1}\sum_{a_1<\cdots<a_n}\prod_{i\le n}a_i^{-p}
 \le\sum_{n\ge1}\frac{\zeta(p)^n}{n!}<\infty.
\]
The limsup set has $d$-dimensional Hausdorff measure zero. Taking $R\uparrow e^B$ and $d\downarrow1/(R-1)$ gives the upper bound for finite $B>\log2$. Let $R\to\infty$ for $B=\infty$ and use the ambient dimension for $B=\log2$. This proves Theorem~\ref{thm:spectrum} for both conventions, as the lower construction tolerates repetitions and the upper bound applies to every greedy expansion.

\section{A single full-dimensional family}
Taking $v_n=\log2$ for all $n$ gives an explicit integer recursion
\[
 M_1=8,\qquad M_{n+1}=(128(n+1))^{3(n+1)}M_n^2.
\]
The compact survivor set for this one schedule already has Hausdorff dimension one. Moreover $2^{-n}\log(1/E_n)$ converges on every surviving branch to the \emph{same} positive finite number
\[
 \log 8+\sum_{k\ge1}2^{-k}\,3(k+1)\log(128(k+1)).
\]
Indeed $\log(1/E_n)=\log M_{n+1}+O(1)$ and the formula follows by unrolling the exact recursion. Also $a_{n+1}/a_n^2\to\infty$ on every branch. The collision count allows substantial freedom in each denominator even with all these requirements imposed simultaneously.

\paragraph{Further questions.} What is the optimal Hausdorff gauge on the full-dimensional set in Theorem~\ref{thm:full}? How does the dimension change if $a_{n+1}/a_n^2$ must remain bounded? Our counting argument takes this ratio to infinity to obtain summable losses. Does every positive algebraic irrational admit eventually greedy best underapproximations, as Erd\H os and Graham suggested? The measure and dimension arguments here do not decide membership for a specified algebraic irrational.

\begin{thebibliography}{9}
\bibitem{kovac} V. Kova\v c, On eventually greedy best underapproximations by Egyptian fractions, \emph{J. Number Theory} \textbf{268} (2025), 39--48. \href{https://doi.org/10.1016/j.jnt.2024.09.004}{doi:10.1016/j.jnt.2024.09.004}.
\bibitem{kt} V. Kova\v c and Q. Tang, Eventually greedy best Egyptian underapproximations of rational numbers via optimal control, 2026. \href{https://arxiv.org/abs/2607.28387}{arXiv:2607.28387}.
\bibitem{nathanson} M. B. Nathanson, Underapproximation by Egyptian fractions, \emph{J. Number Theory} \textbf{242} (2023), 208--234. \href{https://doi.org/10.1016/j.jnt.2022.07.005}{doi:10.1016/j.jnt.2022.07.005}.
\bibitem{wu} J. Wu, On the distribution of denominators in Sylvester expansions (II), \emph{Math. Proc. Camb. Philos. Soc.} \textbf{135} (2003), 421--430. \href{https://doi.org/10.1017/S0305004103006935}{doi:10.1017/S0305004103006935}.
\bibitem{repo} V. Reshetnikov, \emph{ProveIt}, Generalized Riemann Hypothesis article, section on the oriented kernel away from the origin, including the separation lemma, the spectrum theorem and Question \texttt{li:q:decay}. \href{https://github.com/VladimirReshetnikov/ProveIt/blob/ff96fd44ab160e65fb8a9d861243828d07c156cf/Analysis/GeneralizedRiemannHypothesis/Article/sections/li.tex}{Pinned source, revision \texttt{ff96fd44ab}}.
\end{thebibliography}
\end{document}
